\section{Trace Equation Systems}\label{sec:teq}

A loop is described by equations over \emph{traces}, the sequences of values that its variables take, one
value per iteration. This section defines the combinators that the equations use, how a loop is turned into
a system of equations, and what such a system looks like for single loops, for loops whose variables read
each other, and for nested loops.

\subsection{Combinators}
A trace of $n$ steps has $n + 1$ values. Its last value is $\cur(X)$, and the $n$ values before it are the
list $\bl{X}$, the \emph{history} of $X$; the $j$-th earlier value is $\mathsf{earlier}(j, \bl{X})$. The
\emph{counter} $\cnt(i, c, n)$ is the trace $\langle i, i + c, \ldots, i + nc\rangle$ of $n$ steps. Traces
are built and observed by three combinators, each iterating over one or more lists of equal length, the
\emph{dependencies}, whose elements are bound pointwise:
\begin{itemize}[nosep]
\item $\tr{x_0}{D_1, \ldots, D_m}{x, y_1, \ldots, y_m \mapsto u}$, the trace from $x_0$ that steps to $u$ with the
  next values $y_i$ of the $D_i$;
\item $\mapc{D_1, \ldots, D_m}{y_1, \ldots, y_m \mapsto e}$, the list of the values of $e$;
\item $\al{D_1, \ldots, D_m}{y_1, \ldots, y_m \mapsto \varphi}$, $\varphi$ for all positions.
\end{itemize}
In Dafny the three are written \texttt{trace(x0, for D do x, y :: u)}, \texttt{maps(for D yield y :: e)} and
\texttt{invs(for D check y :: φ)}, the history of a trace \texttt{⌊X⌋} and its last value
\texttt{current(X)}. The number of steps of a trace is never a term of the goal: it is the length of the
counter of the iterations, $K = \cnt(0, 1, n)$, and lemmas iterate over counters whose number of steps $n$
they name directly. The type $\mathsf{Nat}$ of such numbers is a datatype with its embedding
$\intf(n)$ into the integers and the inverse $\natf(i)$, which is $0$ for negative $i$.

An update may consist of guarded cases, $\upd{\varphi_1}{e_1}\ \ \upd{\varphi_2}{e_2}$, which take the
result of the case whose guard holds. Where no guard holds the update is $\mathsf{unspecified}$, one
fixed but unknown value per type. The guards of the cases that symbolic execution produces are complete
for the iterations that the loop actually performs, which the iteration constraint (below) states, and
disjoint, since they are the conditions of the branches of the loop body.

A trace may have several accumulators, for variables that read each other:
\[ S = \mathsf{trace}(x_0, y_0,\ \mathsf{for}\ D\ \mathsf{do}\ x, y, d \mapsto u, v) \]
steps both components at once. Its components are read by $\cur(S).0$, $\cur(S).1$ and their histories by
$\bl{S}.0$, $\bl{S}.1$, each an ordinary list that can be a dependency of further combinators. No term of a
tuple type is ever built: a component is only read through these two observers.

\subsection{From a loop to a system}
Symbolic execution of a loop introduces the counter $K = \cnt(0, 1, n)$ of a fresh number $n$ of
iterations and expresses the values of the modified variables after the loop as current values of
traces. It records three kinds of facts:
\begin{itemize}[nosep]
\item the \emph{trace equations}, $X = \tr{x_0}{\bl{K}, \ldots}{x, k, \ldots \mapsto u}$, one per
  trace, whose update $u$ is the effect of one execution of the body;
\item the \emph{iteration constraint}, $\al{\bl{K}, \bl{X}, \ldots}{k, x, \ldots \mapsto \mathit{test}}$:
  the loop test held at the start of every iteration;
\item the \emph{exit}, $\neg\mathit{test}$ for the current values: the loop has stopped.
\end{itemize}
Each variable reads some others in the update. The modified variables are partitioned by the strongly
connected components of this read relation, ordered such that each comes after the components it reads, and
each component becomes \emph{one} trace: a variable in no cycle gets a trace of one accumulator, a cycle
gets a trace of several. Each trace has the counter, the histories of the earlier components it reads, and
the choices of the body among its dependencies. Since the components form an acyclic graph, no trace refers
to a later one, and the system can be read from top to bottom, each equation defining a trace in terms of
the previous ones.

\paragraph{A single loop.}
The linear search \texttt{i := 0; while i < n \&\& a[i] != x \{ i := i + 1; \}} gives
\begin{gather*}
  K = \cnt(0, 1, n') \qquad I = \tr{0}{\bl{K}}{i, k \mapsto i + 1} \qquad i = \cur(I)\\
  \al{\bl{K}, \bl{I}}{k, i \mapsto i < n \wedge a[i] \ne x} \qquad \neg(i < n \wedge a[i] \ne x).
\end{gather*}
The trace of $i$ is a counter, the constraint states the test for the earlier values of $i$, the exit for
the current one.

\paragraph{A cycle.}
In the binary search, \texttt{lo} and \texttt{hi} read each other through the probe \texttt{(lo + hi) / 2},
and the two variables form one trace of two accumulators:
\begin{align*}
  S = \mathsf{trace}(0, n,\ \mathsf{for}\ \bl{K}\ \mathsf{do}\ lo, hi, k \mapsto\ &\upd{G}{(hi + lo)/2 + 1}\ \ \upd{\neg G}{lo},\\
  &\upd{G}{hi}\ \ \upd{\neg G}{(hi + lo)/2})
\end{align*}
with $G$ for $a[(hi + lo)/2] < x$, the exit values $lo = \cur(S).0$, $hi = \cur(S).1$, and the constraint
$\al{\bl{K}, \bl{S}.0, \bl{S}.1}{k, lo, hi \mapsto lo < hi}$. Equations of one accumulator each would be
recursive through each other; one equation over the pair is not. The per-variable view of a component is
still available when a lemma needs it, as a \emph{projection} (Section~\ref{sec:engine:views}).

\paragraph{Nested loops.}
An inner loop is summarized while the body of the outer loop is executed, as a system of its own over a
fresh number of inner iterations. This system is \emph{closed}: its counter, its traces and its exit values
are substituted by their definitions, so that only the number of inner iterations remains fresh, and the
inner loop is a value of the outer update. In the nested sum
\begin{lstlisting}
c := 0; i := 0;
while i < m { j := 0; while j < n { j := j + 1; } c := c + j; i := i + 1; }
\end{lstlisting}
the final value of $j$ in an outer iteration is $\cur(\tr{0}{\bl{\cnt(0, 1, n_7)}}{j, k \mapsto j + 1})$ for
the number $n_7$ of inner iterations of that outer iteration, and the trace of $c$ is
\[ C = \tr{0}{\bl{K}, \seq{N}}{c, k, n_7 \mapsto c + \cur(\tr{0}{\bl{\cnt(0, 1, n_7)}}{j, k' \mapsto j + 1})}. \]
The numbers of inner iterations form a list $\seq{N}$ of the outer loop, a \emph{choice}: a dependency that
nothing defines, about which the constraint of the outer loop states what the inner loop's own constraint
and exit say,
\begin{align*}
  \mathsf{invs}(\mathsf{for}\ \bl{K}, \ldots, \seq{N}\ \mathsf{check}\ k, \ldots, n_7 \mapsto\
    & i < m \wedge \al{\bl{\cnt(0, 1, n_7)}, \ldots}{\ldots \mapsto j < n}\\
    & {} \wedge \neg(\cur(\ldots) < n)).
\end{align*}
Lemmas on single loops apply to the inner loop as it stands in the update, since it is a term like any
other; once the constraint determines the numbers of iterations, e.g.\ $n_7 = \natf(n)$ in every outer
iteration, the choice becomes the constant list $\bl{\mathsf{consts}(n', \natf(n))}$ and the inner traces have
a known number of steps.

\paragraph{Exact, with exits as branch conditions.}
The translation abstracts nothing: the goal is valid iff the program satisfies its assertions, up to the modelled
semantics of the input language, e.g.\ bounded integers with wrap-around for the unsigned types of C. Calls are
inlined from the entry point (\texttt{main} in C), and the body of each loop is checked to execute safely from the
state at the head. Exits by \texttt{break} and \texttt{return} are told apart by branch conditions, not by
assumptions: as assumptions, two exits with different values gave overlapping cases of an enclosing loop, from which
any goal followed, a wrong proof found on SV-COMP.

\paragraph{The goal.}
The goal of a method is that its postcondition follows from the facts of the path through it, with the
traces as variables defined by their equations: $K = \cnt(0, 1, n')$, $S = \mathsf{trace}(\ldots)$, the
exit values $lo = \cur(S).0$, the iteration constraint and the exit. The equations of the outermost loops are
kept as equations, which the engine follows where a lemma needs the definition of a trace variable; inner
loops are closed, as above.
